\section{Introduction}

Particle-based methods are useful for simulating fluids with free surfaces and
changing topology. Smoothed particle hydrodynamics represents a continuum by a
set of moving particles and approximates field quantities using a smoothing
kernel \cite{lucy1977,morris1997}.

The goal of this project is to implement and inspect a compact two-dimensional
SPH solver in C++17. The application includes density and pressure estimation,
viscosity, gravity, boundary particles, and an interactive SFML renderer. This
paper describes the numerical model, the software architecture, and the tests
used to assess the implementation.

\section{Research Questions}

The evaluation focuses on three questions:
\begin{enumerate}
    \item Does the solver preserve the expected qualitative behavior of a small
    fluid scene?
    \item How does the radius-based neighbor search affect the simulation's
    computational work?
    \item Which implementation parameters have the greatest visible effect on
    stability and motion?
\end{enumerate}